% EasyVision paper -- PLAN AND SKELETON (not for submission).
% Red text marks what must be filled in once the EasyVision code runs.
% Everything else is written to be kept.
\documentclass[11pt,a4paper]{article}

\usepackage[T1]{fontenc}
\usepackage[utf8]{inputenc}
\IfFileExists{lmodern.sty}{\usepackage{lmodern}}{\usepackage{mathptmx}}
\usepackage{microtype}
\usepackage[margin=2.5cm]{geometry}
\usepackage{amsmath}
\usepackage{graphicx}
\usepackage{booktabs}
\usepackage{tabularx}
\usepackage{xcolor}
\usepackage[numbers,sort&compress]{natbib}
\usepackage[hidelinks]{hyperref}
\usepackage[capitalise,noabbrev]{cleveref}
\usepackage{authblk}

\newcommand{\EV}{\textsc{EasyVision}}
\newcommand{\EC}{\textsc{EasyClassifier}}
\newcommand{\ER}{\textsc{EasyResearch}}
\newcommand{\TBD}[1]{\textcolor{red}{[#1]}}
\newcommand{\pp}{\,pp}

\title{\EV: How Much Do Evaluation Shortcuts Inflate Image-Classification Results?\\ A Tool That Prevents Them and a Measurement of the Gap}

\author[1]{Ahmad B. A. Hassanat\thanks{Corresponding author: \href{mailto:ahmad.hassanat@gmail.com}{ahmad.hassanat@gmail.com}}}
\author[2]{Ghada A. Altarawneh\thanks{\href{mailto:ghadaa@mutah.edu.jo}{ghadaa@mutah.edu.jo}}}
\affil[1]{Faculty of Information Technology, Mutah University, Karak, Jordan}
\affil[2]{Department of Accounting, Mutah University, Karak, Jordan}
\date{\TBD{Plan and skeleton, 10 October 2026}}

\begin{document}
\maketitle

\begin{abstract}
Image classifiers in applied research are often reported with accuracies that do not hold on new data, and the cause is frequently the evaluation rather than the network: images of the same subject on both sides of the split, augmented copies of test images in the training set, near-duplicates, and the reporting of the best of several models. Such errors are invisible in a paper whose methods section does not say when augmentation was applied or how images were split. We present \EV, an open-source tool for researchers who do not program, which makes these errors impossible by design: augmentation exists only inside each training split, splits follow subjects, copies (exact, flipped, rotated, resized or recoloured) are removed before the analysis, and the selected model receives a separate final score. On \TBD{N} public datasets from \TBD{fields}, we measured the same models under five evaluation regimes. Compared with the honest regime, augmenting before splitting inflated balanced accuracy by \TBD{X}\pp{} on average, image-level instead of subject-level splitting by \TBD{X}\pp, leaving near-duplicates in by \TBD{X}\pp, and adding flipped and rotated copies as new images by \TBD{X}\pp. On random labels the leaky regimes reported up to \TBD{X}\% balanced accuracy, against 50\% for the honest one. \EV{} is available under the MIT license \TBD{PyPI, GitHub, Zenodo DOI}.
\end{abstract}

\noindent\textbf{Keywords:} computer vision; image classification; data leakage; data augmentation; near-duplicates; subject-level splitting; reproducibility; no-code

%=====================================================================
\section{Introduction}
\label{sec:intro}

Deep networks have made image classification accessible to researchers in medicine, biology, agriculture, materials science and the humanities, many of whom do not program and have no training in how such models should be evaluated. Published accuracies are often high, yet many do not survive contact with new data. In a review of models for COVID-19 detection from chest radiographs and CT scans, none of the models examined was judged to be of potential clinical use, largely because of methodological flaws and biases, including duplicated images and patients across sets and datasets whose classes came from different sources \citep{roberts2021covid}. In brain MRI, splitting 2D slices rather than subjects between training and test inflated slice-level test accuracy by 29--55\% across four datasets, and on data with random labels the leaky split still produced about 96\% accuracy against the 50\% expected \citep{yagis2021leakage}. About 3.3\% of the CIFAR-10 test images and 10\% of the CIFAR-100 test images have duplicates in the training set, and well-known networks lost accuracy on a duplicate-free test set \citep{barz2020cifar}. Models also learn the hospital, the scanner or text markers instead of the disease \citep{zech2018pneumonia,geirhos2020shortcut,degrave2021shortcuts}.

Leakage is a problem of machine-learning-based science in general \citep{kapoor2023leakage}, and medical imaging has its own long list of methodological failures \citep{varoquaux2022failures}. What makes the image case particularly hard to detect is that the decisive information is rarely in the paper: a reader cannot tell whether augmentation was applied before or after splitting, whether several images came from one patient, or whether the ``10{,}000 images'' were in fact 1{,}000 images saved ten times in rotated form.

In earlier work we built \EC{} \citep{hassanat2026easyclassifierpaper} for tabular classification and \ER{} \citep{hassanat2026easyresearch} for regression, forecasting and signals, around one principle: the honest procedure is the default and cannot be bypassed by accident. We showed for tables that the common practice overstated balanced accuracy by 2.9 percentage points on average and up to ten on small data. Images offer more ways to inflate a result, and probably larger ones. This paper makes two contributions:
\begin{enumerate}
  \item \EV, a tool that applies the same principle to images (\cref{sec:design}): augmentation only inside each training split, subject-level splitting, a mandatory copy check, early stopping and model selection on training data only, and a report that states every check.
  \item A measurement (\cref{sec:protocol,sec:results}) of how much each common shortcut inflates the reported result, on the same models and data, with a random-label control for each.
\end{enumerate}

%=====================================================================
\section{Related work}
\label{sec:related}

\paragraph{Leakage and pitfalls.} \TBD{Expand: Kapoor \& Narayanan's taxonomy \citep{kapoor2023leakage}; Lones' guide \citep{lones2021pitfalls}; selection bias after model comparison \citep{varma2006bias,cawley2010overfitting}; Varoquaux \& Cheplygina \citep{varoquaux2022failures}.}

\paragraph{Leakage specific to images.} \TBD{Subject-level splitting \citep{yagis2021leakage}; near-duplicates \citep{barz2020cifar}; shortcuts and confounding by site \citep{zech2018pneumonia,geirhos2020shortcut,degrave2021shortcuts}; the COVID-19 review \citep{roberts2021covid}. Add: augmentation-before-split studies, and patch-level leakage in histopathology.}

\paragraph{Augmentation.} Augmentation is a legitimate way to make training harder to overfit \citep{shorten2019augmentation}. \TBD{Add evidence on offline augmentation before splitting and on augmentation used to balance classes.}

\paragraph{Tools.} \TBD{No-code and low-code image tools (e.g.\ general workbenches, AutoML services); what each does about grouping, augmentation scope and copies. Position \EV{} as the first to enforce all of them by default --- check this claim.}

%=====================================================================
\section{Design of \EV}
\label{sec:design}

\EV{} follows the rules of \EC{} and \ER, applied to images:
\begin{enumerate}
  \item Everything learned from data is learned on training images only, inside every split.
  \item Splits follow the subject (patient, specimen, plant, slide, video) whenever images share a source.
  \item Copies are never data: exact, flipped, rotated, resized or recoloured copies are removed before the analysis, and only unique images are counted.
  \item The selected model receives a separate, honest final score (nested cross-validation for small data, an untouched subject-wise test set for larger data).
  \item Default settings are fixed in advance, with fixed seeds; there is no hyperparameter search. Fine-tuning a pretrained network with one documented recipe is training, not tuning.
  \item Every run writes a plain-language report with a Methods paragraph, the leak checks performed, figures, predictions, the trained model and a log.
\end{enumerate}

\begin{table}[t]
\centering\small
\caption{Common leaks in image classification and what \EV{} does about each.}
\label{tab:leaks}
\begin{tabularx}{\linewidth}{@{}>{\raggedright\arraybackslash}p{3.2cm}>{\raggedright\arraybackslash}X>{\raggedright\arraybackslash}X@{}}
\toprule
Leak & How it happens & What \EV{} does \\
\midrule
Augmentation before splitting & the dataset is enlarged offline, then split & augmentation exists only as an on-the-fly transform of training images; augmented images are never written to disk \\
Augmentation to balance classes & the minority class is multiplied before splitting & not offered; imbalance is handled with class weights and balanced accuracy \\
Several images per subject & patients, specimens, slides, slices, frames, tiles & a subject column is part of the data layout; all images of a subject stay in one fold \\
Copies passed off as new images & the same image saved again, flipped, rotated, cropped or recoloured & mandatory copy check before anything else; copies removed and counted \\
Near-duplicates & re-uploads, crops, the same photograph in two classes & perceptual hashing across the dataset; duplicates merged into one subject or removed \\
Normalisation from all images & mean and SD computed on every image & fixed published normalisation of pretrained networks; anything learned is learned on training images \\
Early stopping on the test fold & the best epoch on the test data is kept & early stopping on a validation part of the training images, by subject \\
Best of several networks & the winner's comparison score is reported & nested CV or an untouched subject-wise test set \\
Classes from different sources & each class from another hospital, camera or website & a check whether image properties (size, format, compression, metadata) predict the class \\
\bottomrule
\end{tabularx}
\end{table}

\paragraph{Copy check.} \TBD{Describe the four stages (file and pixel hashes; perceptual hash over the eight flip-and-rotation variants; perceptual-hash distance for resized and recoloured copies; feature similarity for crops, confirmed by the user) and the rule that an image found in two classes stops the analysis. Report its precision and recall on synthetic copies (\cref{tab:copycheck}).}

\paragraph{Models.} \TBD{Linear probe on frozen pretrained features; fine-tuned small ResNet, EfficientNet and vision transformer with one documented recipe; colour-and-texture features with a random forest; KNN with the Hassanat distance \citep{hassanat2014arxiv,hassanat2014distance,hassanat2014distance2} on extracted features. Give the recipe (optimizer, learning rate, epochs, early-stopping rule) and the weights used, with licences.}

\paragraph{Outputs.} \TBD{The standard results folder, plus ``What was checked'' and a reviewer checklist filled in for the analysis.}

%=====================================================================
\section{Experimental protocol}
\label{sec:protocol}

\paragraph{Evaluation regimes.} The same models and data are evaluated five ways:
\begin{enumerate}
  \item[R1] \emph{Honest}: subject-wise splits, augmentation inside training only, copies removed, early stopping on training data, separate final score.
  \item[R2] \emph{Augmentation before splitting}: each image augmented offline $a$ times, then the enlarged set split at random.
  \item[R3] \emph{Image-level splitting}: images split at random, ignoring subjects.
  \item[R4] \emph{Near-duplicates left in}: the copy check switched off.
  \item[R5] \emph{Copies as new images}: flipped and rotated copies added to the dataset as new samples before splitting (the practice of \cref{sec:design}).
\end{enumerate}
Each regime is also run with the class labels permuted at random (by subject), where the true performance is chance.

\paragraph{Truth.} \TBD{For each dataset, an untouched subject-wise external set (or an official external test set where one exists) gives the true performance of the model each regime selects; the inflation of a regime is its reported score minus that truth.}

\paragraph{Datasets.} \TBD{Choose public datasets with subject identifiers, from several fields. Candidates to check for licence, size and identifiers: dermoscopy with lesion identifiers; histopathology with patient identifiers in file names; chest radiographs with patient identifiers; a plant-disease or animal-identification set with individual identifiers; one non-medical set with known near-duplicates. Record for each: images, subjects, classes, copies found by the copy check.}

\begin{table}[t]
\centering\small
\caption{Datasets. \TBD{to be filled}}
\label{tab:datasets}
\begin{tabular}{@{}llrrrr@{}}
\toprule
Dataset & Field & Images & Subjects & Classes & Copies found \\
\midrule
\TBD{dataset 1} & \TBD{} & -- & -- & -- & -- \\
\TBD{dataset 2} & \TBD{} & -- & -- & -- & -- \\
\TBD{dataset 3} & \TBD{} & -- & -- & -- & -- \\
\TBD{dataset 4} & \TBD{} & -- & -- & -- & -- \\
\bottomrule
\end{tabular}
\end{table}

\paragraph{Measures and statistics.} Balanced accuracy is the primary measure, per image and, where subjects have several images, per subject (averaged predictions). Inflation is compared between regimes with paired Wilcoxon signed-rank tests over datasets and seeds, with a sensitivity analysis that averages seeds within each dataset first. \TBD{Number of seeds; computing budget; GPU or CPU.}

%=====================================================================
\section{Results}
\label{sec:results}

\begin{table}[t]
\centering\small
\caption{Reported balanced accuracy (\%) under each regime and the true performance on the untouched external set. \TBD{to be filled}}
\label{tab:results}
\begin{tabular}{@{}lrrrrrr@{}}
\toprule
Dataset & R1 honest & R2 aug.\ first & R3 image split & R4 near-dup. & R5 copies & Truth \\
\midrule
\TBD{dataset 1} & -- & -- & -- & -- & -- & -- \\
\TBD{dataset 2} & -- & -- & -- & -- & -- & -- \\
\TBD{dataset 3} & -- & -- & -- & -- & -- & -- \\
\TBD{dataset 4} & -- & -- & -- & -- & -- & -- \\
\midrule
Random labels & -- & -- & -- & -- & -- & 50.0 \\
\bottomrule
\end{tabular}
\end{table}

\begin{table}[t]
\centering\small
\caption{Copy check on synthetic copies of known kind. \TBD{to be filled}}
\label{tab:copycheck}
\begin{tabular}{@{}lrr@{}}
\toprule
Kind of copy & Found (recall) & False alarms (precision) \\
\midrule
Exact file / same pixels & -- & -- \\
Flipped or rotated by 90/180/270$^\circ$ & -- & -- \\
Resized, recompressed or recoloured & -- & -- \\
Crops and small rotations (user-confirmed) & -- & -- \\
\bottomrule
\end{tabular}
\end{table}

\TBD{Figure: inflation per regime and dataset (as Fig.~4 of the \EC{} paper). Figure: examples of copies found. Figure: occlusion maps for one shortcut case.}

%=====================================================================
\section{Discussion}
\label{sec:discussion}

\TBD{Which shortcut inflates most, and on which kind of data; why random labels are the clearest demonstration; what a reviewer should ask for (the checklist); what \EV{} cannot check (pretraining overlap with public benchmarks, shortcuts that are biologically plausible).}

\paragraph{Limitations.} \TBD{Classification only in phase 1 (segmentation and detection later); 2D images (3D volumes analysed as slices with subject-level evaluation); computing time of fine-tuning on a CPU; the copy check cannot prove the absence of copies.}

\section*{Acknowledgements}
\TBD{As in the \EC{} paper, including the AI-assistance disclosure.}

\section*{Software and data availability}
\TBD{PyPI, GitHub, Zenodo DOI; benchmark code and results.}

\bibliographystyle{plainnat}
\bibliography{references}

\end{document}
